% TOP_PROBLEM: {"id": 3194, "title": "Coloring the union of degenerate graphs", "category_id": 3, "category": {"id": 3, "name": "graph_theory", "display_name": "Graph Theory", "description": "Problems involving graphs, networks, and their properties.", "slug": "graph-theory", "order_index": 3, "created_at": "2026-07-31T15:26:25.671Z"}, "status": "open", "problem_number": "OPG-46533", "set_id": 12}
% FIRST_POSTED: 2026-10-01
% ENTRY_KIND: statement_only
% UnsolvedMath ID 3194; source catalogue code OPG-46533
% !TeX program = lualatex
% Definition and sources only; this file is not an LLM solution attempt.
\documentclass[11pt,a4paper]{article}
\usepackage[margin=25mm]{geometry}
\usepackage{fontspec}
\setmainfont{lmroman10-regular.otf}[BoldFont=lmroman10-bold.otf,
 ItalicFont=lmroman10-italic.otf,BoldItalicFont=lmroman10-bolditalic.otf]
\usepackage{amsmath,amssymb,mathtools}
\usepackage{unicode-math}
\setmathfont{latinmodern-math.otf}
\AtBeginDocument{\let\setminus\smallsetminus}
\usepackage{xurl,hyperref}
\hypersetup{colorlinks=true,urlcolor=blue}
\setcounter{secnumdepth}{0}
\setlength{\parindent}{0pt}
\setlength{\parskip}{5pt}
\setlength{\emergencystretch}{3em}
\begin{document}
\section*{Coloring the union of degenerate graphs}\label{3194}
{\small UnsolvedMath ID 3194 · OPG-46533 · Graph Theory · OpenGarden}
\subsection{Definitions and mathematical statement}
A finite simple graph $H$ is $d$-degenerate if every
nonempty subgraph of $H$ has a vertex of degree at
most $d$. Equivalently, its vertices can be ordered
so that each has at most $d$ neighbours later in
the order. A $1$-degenerate graph is precisely a
forest, that is, a graph with no cycle.

Let $G_1=(V,E_1)$ be $1$-degenerate and
$G_2=(V,E_2)$ be $2$-degenerate, on a common finite
vertex set $V$. If originally their vertex sets
differ, enlarge each by isolated vertices first.
Their union is the simple graph
$G=(V,E_1\cup E_2)$; overlapping edges are counted
only once.

The conjecture asserts that every such union has
a proper vertex colouring with at most five colours:
\[
 \chi(G_1\cup G_2)\leq5.
\]
Equivalently, there is $c:V\to\{1,2,3,4,5\}$
with $c(u)\neq c(v)$ whenever $uv$ belongs to
either edge set. The decomposition is part of
the hypothesis, but its two degeneracy orderings
need not agree.

The elementary product construction uses six
colours: greedily two-colour the forest, three-colour
the $2$-degenerate graph, and label each vertex
by its pair of colours. Every edge has different
labels in at least one coordinate. The problem
asks whether this universal bound can always be
reduced to five, while still simultaneously
respecting all edges of both graphs.
\subsection{Short English statement}
Is the union of a forest and a 2-degenerate graph always properly colourable with five colours?
\subsection{Sources}
\begin{itemize}
\item[S1] ULAM.AI and the UnsolvedMath Contributors, supplied problems.json, record 3194 (OPG-46533), OpenGarden collection. Catalogue identity and source transcription; CC BY 4.0.
\url{https://huggingface.co/datasets/ulamai/UnsolvedMath}.
\item[S2] Open Problem Garden, Coloring the union of degenerate graphs. Original problem and discussion; the mathematical formulation below is expanded from the supplied source record.
\url{http://www.openproblemgarden.org/op/coloring_the_union_of_degenerate_graphs}.
\end{itemize}
\end{document}
